\section{SISC: patterns on real data}

\begin{frame}{S\&P 500: learned motifs ($K=14$)}
\centering
\includegraphics[width=0.88\textwidth]{figures/sisc/sp500/k14/patterns.pdf}
\note{Motifs learned by SISC with K=14. Visually: rebound, plateau, accelerated drawdown, and a clear head-and-shoulders shape around motifs 7--9. K=14 is the setting used in the paper. Centroid lengths span $[\lmin,\lmax]=[10,21]$ trading days.}
\end{frame}

\begin{frame}{GOOG: learned motifs ($K=11$)}
\centering
\includegraphics[width=0.88\textwidth]{figures/sisc/goog/k11/patterns.pdf}
\note{On GOOG we drop to K=11 because the post-SISC segment count shrinks proportionally to the underlying series length. The library still recovers the main shapes: trending up/down, V/inverted-V, and a couple of consolidation motifs.}
\end{frame}

\begin{frame}{ZC$=$F (corn futures): learned motifs ($K=11$)}
\centering
\includegraphics[width=0.88\textwidth]{figures/sisc/zcf/k11/patterns.pdf}
\note{Corn futures: regime is dominated by mean-reverting motifs and sharp impulses around delivery cycles. The motif library is qualitatively different from the index, which is reflected in the downstream MAPE.}
\end{frame}
